\begin{proof}
\textbf{Degenerate case.} If \(\displaystyle\sum_{i=1}^{n}b_i^{2}=0\), then each \(b_i=0\). Hence
\[
\sum_{i=1}^{n}a_i b_i =0,
\qquad 
\Bigl(\sum_{i=1}^{n}a_i^{2}\Bigr)\Bigl(\sum_{i=1}^{n}b_i^{2}\Bigr)=0,
\]
so \(0^{2}=0\le 0\). The equality condition is satisfied because the zero vector is proportional to any vector.

\medskip

\textbf{Non‑degenerate case.} Assume
\[
A:=\sum_{i=1}^{n}b_i^{2}>0
\]
and define, for a real parameter \(t\),
\[
Q(t):=\sum_{i=1}^{n}\bigl(a_i-t b_i\bigr)^{2}.
\]
Each summand is non‑negative, therefore
\begin{equation}
Q(t)\ge 0\qquad\text{for all }t\in\mathbb{R}. \label{eq:Qnonneg}
\end{equation}

Expanding \(Q(t)\) and using \((x-y)^{2}=x^{2}-2xy+y^{2}\) gives
\begin{align*}
Q(t)
&=\sum_{i=1}^{n}\bigl(a_i^{2}-2t a_i b_i+t^{2}b_i^{2}\bigr) \\
&= \underbrace{\sum_{i=1}^{n}a_i^{2}}_{C}
   -2t\underbrace{\sum_{i=1}^{n}a_i b_i}_{S_{ab}}
   +t^{2}\underbrace{\sum_{i=1}^{n}b_i^{2}}_{A}.
\end{align*}
Thus \(Q(t)=A t^{2}+B t +C\) with
\[
A=\sum_{i=1}^{n}b_i^{2}>0,\qquad
B=-2S_{ab},\qquad
C=\sum_{i=1}^{n}a_i^{2}.
\]

A quadratic polynomial with leading coefficient \(A>0\) that satisfies
\(At^{2}+Bt+C\ge0\) for every real \(t\) must have a non‑positive
discriminant:
\[
B^{2}-4AC\le 0.
\]
Applying this to \eqref{eq:Qnonneg} we obtain
\[
(-2S_{ab})^{2}-4\Bigl(\sum_{i=1}^{n}b_i^{2}\Bigr)\Bigl(\sum_{i=1}^{n}a_i^{2}\Bigr)\le0.
\]
Dividing by \(4>0\) yields
\[
S_{ab}^{2}\le\Bigl(\sum_{i=1}^{n}a_i^{2}\Bigr)\Bigl(\sum_{i=1}^{n}b_i^{2}\Bigr),
\]
i.e.
\[
\Bigl(\sum_{i=1}^{n}a_i b_i\Bigr)^{2}
\le
\Bigl(\sum_{i=1}^{n}a_i^{2}\Bigr)
\Bigl(\sum_{i=1}^{n}b_i^{2}\Bigr).
\]
This is the Cauchy–Schwarz inequality.

\medskip

\textbf{Equality condition.} Equality occurs precisely when the
discriminant is zero, i.e. \(B^{2}-4AC=0\). Then the quadratic has the
unique real root
\[
t_{0}= -\frac{B}{2A}= \frac{S_{ab}}{\sum_{i=1}^{n}b_i^{2}}.
\]
Since \(Q(t_{0})=0\) and each term \((a_i-t_{0}b_i)^{2}\) is
non‑negative, we must have \(a_i-t_{0}b_i=0\) for every \(i\); thus
\(a_i=\lambda b_i\) with \(\lambda:=t_{0}\). Conversely, if such a scalar
\(\lambda\) exists, then every term \((a_i-\lambda b_i)^{2}\) vanishes,
so \(Q(\lambda)=0\) and the discriminant is zero, giving equality in the
Cauchy–Schwarz inequality. This also covers the degenerate case where
all \(b_i=0\) (or all \(a_i=0\)), because any \(\lambda\) satisfies
\(a_i=\lambda b_i\).

Hence, for all real families \((a_i)_{i=1}^{n}\) and \((b_i)_{i=1}^{n}\),
\[
\boxed{\displaystyle
\Bigl(\sum_{i=1}^{n}a_i b_i\Bigr)^{2}
\le
\Bigl(\sum_{i=1}^{n}a_i^{2}\Bigr)
\Bigl(\sum_{i=1}^{n}b_i^{2}\Bigr)}
\]
with equality iff there exists \(\lambda\in\mathbb{R}\) such that
\(a_i=\lambda b_i\) for every \(i\) (including the case where one of
the vectors is the zero vector). \(\square\)
\end{proof}